\section*{अध्याय \ref{ch:g4:clouds-rain-snow} --- बादल, बारिश और हिम}

\begin{solution}{exo:g4:clouds-rain-snow:1}
\omterm{def:g4:clouds-rain-snow:cloud}{बादल} पानी की बहुत-सी नन्ही बूँदों, या बर्फ़ के बहुत-से नन्हे टुकड़ों से बना होता
है, जो हवा में तैर रहे हैं। \omterm{def:g4:clouds-rain-snow:cloud}{कोहरा} ऐसा \omterm{def:g4:clouds-rain-snow:cloud}{बादल} है जो ज़मीन को छूता है।
\end{solution}

\begin{solution}{exo:g4:clouds-rain-snow:2}
(क) \omterm{def:g4:clouds-rain-snow:cloud-families}{स्तरी मेघ}। (ख) \omterm{def:g4:clouds-rain-snow:cloud-families}{पक्षाभ मेघ}। (ग) \omterm{def:g4:clouds-rain-snow:cloud-families}{कपासी मेघ}। (घ) \omterm{def:g4:clouds-rain-snow:cloud-families}{कपासी वर्षी मेघ}।
\end{solution}

\begin{solution}{exo:g4:clouds-rain-snow:3}
(क) \omterm{def:g4:clouds-rain-snow:precipitation}{ओले}। (ख) \omterm{def:g4:clouds-rain-snow:precipitation}{बारिश}। (ग) \omterm{def:g4:clouds-rain-snow:precipitation}{हिम}।
\end{solution}

\begin{solution}{exo:g4:clouds-rain-snow:4}
(क) \omterm{def:g4:clouds-rain-snow:rain-gauge}{वर्षामापी}। (ख) तापमापी। (ग) \omterm{def:g4:clouds-rain-snow:wind-vane}{वायु-दिशासूचक}।
\end{solution}

\begin{solution}{exo:g4:clouds-rain-snow:5}
17 मिलीमीटर \omterm{def:g4:clouds-rain-snow:precipitation}{बारिश} हुई है। $20 - 17 = 3$: 3 मिलीमीटर और \omterm{def:g4:clouds-rain-snow:precipitation}{बारिश} पानी को 20 के
निशान तक पहुँचा देगी।
\end{solution}

\begin{solution}{exo:g4:clouds-rain-snow:6}
सबसे गीला दिन बुधवार था (12 मिलीमीटर)। सोमवार और शुक्रवार को \omterm{def:g4:clouds-rain-snow:precipitation}{बारिश} नहीं हुई।
गुरुवार को 7 मिलीमीटर \omterm{def:g4:clouds-rain-snow:precipitation}{बारिश} हुई।
\end{solution}

\begin{solution}{exo:g4:clouds-rain-snow:7}
\omterm{def:g4:clouds-rain-snow:precipitation}{हिम} के फाहे ठंडे \omterm{def:g4:clouds-rain-snow:cloud}{बादल} में ऊँचाई पर बनते हैं। ऊपर \omterm{def:g1:landscapes-around-us:landform}{पहाड़ी} गाँव में नीचे तक पूरी
हवा जमा देने वाली ठंडी है, इसलिए वे \omterm{def:g4:clouds-rain-snow:precipitation}{हिम} बनकर पहुँचते हैं। और नीचे, \omterm{def:g1:landscapes-around-us:landform}{घाटी} में,
हवा \SI{0}{\celsius} से गरम है: उससे होकर नीचे आते हुए फाहे पिघल जाते हैं, और
ज़मीन तक \omterm{def:g4:clouds-rain-snow:precipitation}{बारिश} बनकर पहुँचते हैं।
\end{solution}

\begin{solution}{exo:g4:clouds-rain-snow:8}
दंड आलेख में 3, 0, 1, 9, 6, 0 और 1 ऊँचाई के सात दंड हैं (मंगलवार और शनिवार को
कोई दंड नहीं)। हफ़्ते में $3 + 0 + 1 + 9 + 6 + 0 + 1 = 20$ मिलीमीटर \omterm{def:g4:clouds-rain-snow:precipitation}{बारिश} हुई।
गुरुवार ने 20 में से 9 मिलीमीटर दिए: हफ़्ते की \omterm{def:g4:clouds-rain-snow:precipitation}{बारिश} का $\frac{9}{20}$।
\end{solution}

\begin{solution}{exo:g4:clouds-rain-snow:9}
दीवारें और पेड़ अपने पास की ज़मीन को \omterm{def:g4:clouds-rain-snow:precipitation}{बारिश} से बचा लेते हैं, ख़ासकर जब हवा \omterm{def:g4:clouds-rain-snow:precipitation}{बारिश}
को तिरछा उड़ाती है, और पत्तियों से बूँदें टपककर \omterm{def:g4:clouds-rain-snow:rain-gauge}{वर्षामापी} में गिर सकती हैं: उनके
पास \omterm{def:g4:clouds-rain-snow:rain-gauge}{वर्षामापी} बहुत कम, या बहुत ज़्यादा पानी पकड़ेगा। तापमापी को हवा का तापमान
नापना है; सीधी धूप में ख़ुद धूप उसे गरम कर देगी और वह बहुत ज़्यादा पढ़ेगा। सफ़ेद
झिलमिलीदार बक्सा धूप को रोकता है और हवा को अंदर आने देता है।
\end{solution}

\begin{solution}{exo:g4:clouds-rain-snow:10}
\omterm{def:g4:clouds-rain-snow:wind-vane}{वायु-दिशासूचक} उधर इशारा करता है जिधर से हवा आती है: हवा पश्चिम से आ रही है।
वह पूर्व की ओर जा रही है।
\end{solution}

\begin{solution}{exo:g4:clouds-rain-snow:11}
एक वर्ग मीटर पर एक मिलीमीटर \omterm{def:g4:clouds-rain-snow:precipitation}{बारिश} एक लीटर है। एक वर्ग मीटर पर 8 मिलीमीटर
\omterm{def:g4:clouds-rain-snow:precipitation}{बारिश} 8 लीटर देती है। 500 वर्ग मीटर पर: $500 \times 8 = 4000$ लीटर। 10 लीटर
की झारियों से: $4000 \div 10 = 400$ झारियाँ।
\end{solution}
